
\subsection{Pinch structure and scalar cross-checks}
\label{sec:pinch}

Feynman parameter space provides a complementary
description of the DPS pinch and its local power
counting. We verify the scalar polynomial identities
and local pinch powers independently of the fitted
amplitude coefficients. The candidate scalings are
taken from a separate momentum-space survey; the
checks here do not repeat that enumeration or replace
the regulated region integrations of
Secs.~\ref{app:regions} and~\ref{sec:r0regions}.

\paragraph{The one-loop pinch: }
For $p_3=xp_1+yp_2+p_{3\perp}$, with
$p_{3\perp}^2=-p_T^2$ and $0<x,y<1$,
define $\lambda^2=p_T^2/s$.
The four daughter propagators of the scalar
box have Feynman parameters
$\alpha_k,\alpha_{\bar k},\alpha_r,
\alpha_{\bar r}$.
The second Symanzik polynomial is exactly
\begin{equation}
\frac{F}{s}=-f_1f_2+\lambda^2(\alpha_k+\alpha_{\bar k})(\alpha_r+\alpha_{\bar r}),
\qquad
f_1=x\alpha_k-(1-x)\alpha_{\bar k},\quad
f_2=y\alpha_r-(1-y)\alpha_{\bar r}.
\label{eq:pinchparameter}
\end{equation}
%
At $p_T=0$, both $F$ and its gradient
vanish on $f_1=f_2=0$ within the
positive Feynman parameter domain.
These conditions are the two light-cone
components of the Landau loop equation.
They correspond to two hard vertices
coinciding at the DPS pinch~\cite{Coleman:1965xm,Gaunt:2011xd}.
The singularity is therefore an
interior pinch rather than an endpoint
singularity in the original parameters.
%
In massless $2\to2$ scattering,
Ref.~\cite{Gardi:2024dissecting}
identified analogous parameter-space
pinches beginning at three loops
in particular nonplanar graphs.
With two massive final-state particles,
the DPS box develops such a pinch
already at one loop.
It arises from cancellations between
terms of opposite sign in $F$,
associated with two distinct
hard subgraphs.
%
Taking the normal coordinates
$f_1,f_2\sim\lambda$ gives a
measure proportional to $\lambda^2$,
while $F^{-2}\sim\lambda^{-4}$.
The scalar box therefore scales
as $\lambda^{-2}$ in four dimensions.
The QCD splitting numerators provide
two powers of transverse momentum,
changing this to $\lambda^0$,
up to logarithms.
The scalar singularity alone does not
determine the power of the QCD amplitude.

\paragraph{Input from momentum space: }
The candidate scalings come from a
separate survey of 54 two-loop graphs
containing the DPS box, together with
16 additional loop-momentum
parametrisations.
%
The box momentum is assigned its
DPS pinch or a hard, collinear, soft
or intermediate scaling.
The additional loop momentum is
assigned hard, collinear, soft,
Glauber, ultrasoft or intermediate
longitudinal scalings.
%
Each record specifies the scaling
$|p_e^2|\sim\lambda^{w_e}s$
of the squared momentum of each
internal line $e$, the momentum
routing and the type of reduced
diagram. The supplied labels
distinguish DPS-type configurations,
with two hard vertices, from
single parton scattering (SPS)
configurations, with one hard
vertex producing both final-state
particles.
%
The pointwise power $P$ includes
the integration measure, propagator
denominators and QCD numerators,
evaluated at generic points of the
assigned scaling. The survey uses
Feynman gauge and the axial gauges
$n=p_1+p_2$ and $n=3p_1+2p_2$.
The scalar power $P_{\rm den}$
instead sets the numerators to unity.

The released \texttt{vectors.json.xz}
contains the squared-momentum scalings,
momentum assignments and original
surface labels. The polynomial
checks below are reproducible from
these inputs, but the programs used
to generate the original momentum-space
survey are not included.
Consequently, neither the candidate
scalings nor their DPS/SPS labels
are independently determined by the
parameter-space calculation.

\paragraph{Two-loop polynomials and local powers: }
For an $L$-loop scalar graph with
$E$ propagators, the Feynman parameter
integrand, apart from the projective
measure and overall factors, is
$U^a/(F-i0)^b$, where
$a=E-2(L+1)$ and $b=E-2L$
in four dimensions.
These expressions describe the
integrand without assuming that
the scalar integral is finite.
%
We construct the first and second
Symanzik polynomials from spanning
trees and two-forests and independently
from the determinant and adjugate
of the loop-momentum quadratic form.
The two constructions agree
identically in $x,y,\lambda^2$
for all 54 canonical two-loop
graphs containing the DPS box.
%
For a supplied parameter scaling
$\alpha_e\sim\lambda^{-w_e}$,
write
$U\sim\lambda^{\mu_U}U_0$
and $F\sim\lambda^{\mu_F}F_0$.

The survey contains 2823 records,
representing 2363 distinct
graph-parameter vectors in
70 parametrisations.
Of these, 1709 records, corresponding
to 1305 distinct vectors, have
a leading polynomial $F_0$ with
exactly two factors that change
sign within the positive parameter
domain. The factors are associated
with the two incoming parton families.
%
This algebraic property does not
by itself identify the reduced
diagram. In particular, the same
two-factor structure occurs in
678 records labelled SPS.
The parameter-space check verifies
the factors without using the
supplied DPS/SPS labels.
%
For the 142 selected DPS-type
scalings with the box at its
DPS pinch and the additional
loop hard, collinear or ultrasoft,
the local scalar powers agree
with the supplied $P_{\rm den}$.
Of these, 127 have
$P_{\rm den}\le0$; the remaining
15 have $P_{\rm den}=2$.
This agreement concerns the
measure and propagators.
It does not test the
numerator-inclusive power $P$.

To determine the local power
near an interior pinch, consider
two normal coordinates $z_1,z_2$.
By multiaffinity, the polynomial
takes the form
$F=Az_1z_2+Bz_1+Cz_2+D$,
with the remaining parameters
held positive.
Its stationary point is
$z_1=-C/A$, $z_2=-B/A$,
and the critical value is
$F_{\rm crit}=(AD-BC)/A$.
%
Suppose the first nonzero term
in $F_{\rm crit}$ occurs at
relative order $\lambda^\delta$
with respect to $\lambda^{\mu_F}$.
For $b\ge2$, the resulting
local scalar power is
\begin{equation}
\omega=-\sum_e w_e+a\mu_U-b\mu_F+\delta(1-b).
\label{eq:pinchlocalpower}
\end{equation}
For $b=1$, the normal integration
is logarithmic.
%
This local analysis applies away
from zeros of the critical value
along the tangent directions and
away from tangent boundaries.
Such configurations require
additional analysis.

A complete region classification
also requires dissecting the
parameter domain at the pinch,
turning the interior singularity
into endpoint singularities.
The regions can then be identified
from the Newton polytopes of
the dissected integrals~\cite{Gardi:2024dissecting}.
The present calculation determines
the stated local powers without
performing this full dissection.

\paragraph{Interjet exchange and gauge dependence: }
For the 16 graphs in which the
additional gluon connects the two
incoming jets, we examine the
common squared-momentum assignments
at the Glauber scaling.
The leading polynomial has no
stationary point with all Feynman
parameters strictly positive.
%
This follows from exact algebra
at arbitrary interior $x,y$:
either a derivative has a fixed
nonzero sign, or every solution
of a necessary stationary equation
forces another derivative to have
a fixed nonzero sign.
%
The result excludes a stationary
point in the positive interior
for these scalar configurations.
It does not exclude boundary
pinches or imply that the regulated
Glauber contribution vanishes.
In particular, it does not
eliminate the Glauber phase
found in Sec.~\ref{app:regions}.

In other graphs, a Glauber
assignment can give the same
squared-momentum scalings as a
collinear or soft assignment.
These possibilities cannot be
distinguished by the scalar
polynomial check alone.
%
The momentum-space survey also
includes intermediate longitudinal
scalings between the soft,
collinear and Glauber regions.
For $gg\to HH$, $gg\to AA$
and $gg\to ZZ$, the only
negative pointwise power found
at these intermediate scalings
is $P=-\tfrac12$, occurring
in individual daughter
self-energy graphs in the
two axial gauges.
No negative power is found
in Feynman gauge at these
scalings.
%
The half-integer powers arise
from gauge-vector terms in
the exchanged-gluon numerator.
They characterise individual
integrands rather than a
demonstrated singularity of
the gauge-invariant amplitude.
Their cancellation in the
complete axial-gauge sum,
including axial-pole
prescriptions and regulated
overlaps, is not established
by the pointwise survey.

\paragraph{Reproduction and scope: }
The released package
\texttt{checks/pinch\_regions/}
contains the supplied scaling
vectors, exact Symanzik
polynomial identities, algebraic
sign checks and regression tests.
%
The scalar polynomial and
local-power calculations are
reproducible from these data.
The original momentum-space
enumeration that selected the
candidate vectors is not
regenerated by this package.
%
The checks establish the
one-loop interior pinch,
the stated two-loop polynomial
identities, the local scalar
powers and the absence of
positive-interior stationary
points in the selected
interjet-exchange configurations.
They provide complementary
evidence for the DPS singular
structure alongside the
two-loop amplitude comparisons
and regulated momentum-region
calculations.
%
A complete classification of
the momentum regions and
renormalised matching of
all overlaps are separate
extensions of these results.
